% Transcription — fonds Alexandre Grothendieck, shelfmark 128, batch 2 (pages 21-31).
% Established from the facsimile archives/batches/128/batch-02.pdf, prepared
% from a copy of 128.pdf supplied by the user (PDF page k = archivists' page
% 20+k, no cover sheet), per the skill transcribe-grothendieck. The folder is
% thirty-one pages; this is the second and last batch. Batch 1 (pages 1-20)
% is transcribed separately and in parallel by another pass and may not
% exist yet, so the notation below was fixed from these pages and from
% folder 126, the first folder of the same « Descentes » group.
%
% Pass: Opus 5.5 (claude-opus-5-5), 2026-09-25 — first pass, unchecked against the pages by a human.
%
% Dating: the folder is « s.d. » in src/content/catalogue.ts; its inventory
% group « Descentes » (cotes 126 to 133) is dated [avant 1970], and \dating
% says both, per the skill's group rule. Nothing on these leaves dates them.
%
% Pages skipped: 30, a blank tan sheet (presumably the back cover of the
% run), nothing in anyone's hand.
%
% Pages summarised: 31. It is a single sheet in another hand — neat, in
% English (« "equivalence ideal" in A' ⊗_k A' », « be a set of generators »,
% « This is still set of generators », « Next I ⊂ Δ ») — on a square-zero
% algebra A' = k + V = k[X_1..X_n]/(X_i)^2 and a cocycle condition
% p13(I) ⊂ p12(I) + p23(I) on an ideal in A' ⊗ A'. Given the folder title
% (« Cf Notes Murre-Levelt ») it is presumably a leaf of those notes, but
% nothing on it names its author, and the file does not name one. It is
% other people's text: one French \note, no transcription. No mark in
% Grothendieck's hand was found on it.
%
% Pages transcribed: 21-29, all his, blue ink on white paper, fast drafting.
% The statements and formulas carry; much connective prose does not and is
% left \ill{}. His pagination: none visible, except the circled letters
% (a), (b), (c) on pp. 26, 27, 29, which number the parts of the proof of
% the Théorème of p. 25 and are kept as he wrote them.
%
% What the batch is about. Descent along a finite algebra A -> A' through
% infinitesimal neighbourhoods, with the non-abelian H^1 of the Amitsur
% (Čech) complex. Pages 21-22: Corollaires 2 and 3 of a « lemme 1 » that is
% presumably at the end of batch 1 (not seen here): for the matrix sheaf
% M_n and Gl_n = M_n^*, and more generally for End(E), Aut(E) with E free,
% the kernels of H^1(A'/A, -) -> H^1(A'_0/A_0, -) correspond bijectively,
% under a hypothesis « End = Aut + Im End_B » of which he gives a footnote;
% hence H^1(A'/A, Gl_n) = {e} iff H^1(A'/A, G_a) = {e} when the reduction
% is trivial. Pages 23-24: Lemme 2 and a variant, the local-algebra lemma
% that makes that hypothesis hold (infinite residue fields, radicial
% residue extensions: every matrix over C can be made invertible by adding
% one coming from B), with an N.B. on a transcendental residue extension
% k(t). Page 25: the Théorème — A semi-local noetherian with radical m,
% A' finite over A, A_n = A/m^{n+1}, A'_n = A' ⊗ A_n: a) equivalence of
% (i) H^1(A'_n/A_n) = 0 for all n, (ii_1), (ii_m) H^1(A'_n/A_n, Gl) trivial,
% (iii) Spec A'_n -> Spec B^(n), B^(n) = H^0(A'_n/A_n), of effective
% descent for flat quasi-coherent modules; b) consequences at the limit
% for A'/A, with B = H^0(A'/A) and modules flat of finite type; c) A
% artinian: X' -> X of effective descent for flat modules. Pages 26-29:
% the proof, by induction on n and the lemmas, then passage to the limit
% by Mittag-Leffler (« M.L. ») on C^.(A'/A) = lim C^.(A'_n/A_n). The last
% line of p. 29 is an isolated A -> A' ⇉ A''.
%
% Notation fixed once for the batch: \mathbb{M}_n for his double-struck M
% (the matrix sheaf), \mathrm{Gl}_n, \mathrm{Gl}(n, C_0); \underline{\mathrm{Aut}},
% \underline{\mathrm{End}}, \underline{\mathrm{Hom}} where he underlines;
% C^{\cdot}(A'/A, -) for his (doubled) C of cochains, C^{\cdot}_D and
% H^0_D where he puts a D under them (descent data); \varprojlim for his
% lim with a left arrow under it. Words he underlines are \emph.
%
% Readings to check first: p. 21 the symbol read Z (Z C^.(A'/A, M_n)) and
% the end of the proof; p. 22 the exponents (N) and N of the bracketed
% term, overwritten; p. 23 the proof of Lemme 2 (the theorem invoked);
% p. 24 nearly all of it; the diagonal N.B. in the margin of p. 25 and
% the fraction « qu. cohér. / plats »; p. 26 middle; p. 27 the bracket.

\documentclass[11pt,a4paper]{article}
% Le préambule commun à toutes les transcriptions du fonds Grothendieck.
%
% Il tient en une idée : **l'incertitude se compose, elle ne se résout pas.**
% Une transcription qui lisse un mot douteux a détruit la seule chose pour
% laquelle on la faisait — un lecteur doit pouvoir savoir, à chaque endroit,
% ce qui était sur le feuillet et ce qui a été ajouté. Les six macros qui
% suivent sont donc l'appareil critique, et non de la décoration.
%
% Les mêmes commandes sont reconnues par `scripts/render.mjs`, qui produit la
% vue de lecture du site. Une macro ajoutée ici doit l'être aussi là-bas,
% faute de quoi le rendu échouera bruyamment — ce qui est voulu : un rendu
% silencieusement faux est pire qu'un rendu absent.

\ProvidesPackage{grothendieck}[2026/08/08 Transcriptions du fonds Grothendieck]

\usepackage{amsmath,amssymb,amsthm}
\usepackage{mathrsfs}
\usepackage{stmaryrd}
% Les diagrammes commutatifs sont partout dans ce fonds : c'est la langue
% même de la géométrie algébrique de Grothendieck, et une transcription qui
% les rendrait en prose ne transcrirait rien.
\usepackage{tikz-cd}
% Français par défaut — c'est la langue des feuillets — mais l'anglais est
% chargé aussi : la traduction et le résumé sont des documents anglais, et la
% typographie française (espaces avant les deux-points) y serait une faute.
% Ces documents basculent par \selectlanguage{english} en tête de corps.
\usepackage[english,french]{babel}
\usepackage{fontspec}
\usepackage{geometry}
\geometry{a4paper,margin=25mm,marginparwidth=28mm}
\usepackage{marginnote}
\usepackage{xcolor}
% ulem, not soul. `soul`'s \st and \ul are built on a character-by-character
% scan that XeLaTeX's font machinery breaks: the document fails with an
% undefined control sequence rather than degrading. ulem does the same two jobs
% and is Unicode-engine safe. `normalem` stops it hijacking \emph, which the
% transcriptions use for Grothendieck's own emphasis.
\usepackage[normalem]{ulem}
\usepackage{hyperref}
\hypersetup{colorlinks=true,linkcolor=black,urlcolor=black,pdfborder={0 0 0}}

% --- Métadonnées du lot -----------------------------------------------------
% Reprises telles quelles par le script de rendu, qui les affiche en tête de
% la vue de lecture. Elles fixent ce que le fichier prétend couvrir : sans
% elles, une transcription flotte sans point d'attache dans un fonds de seize
% mille feuillets.

\newcommand{\folder}[1]{\gdef\grothFolder{#1}}
\newcommand{\batch}[1]{\gdef\grothBatch{#1}}
\newcommand{\leaves}[2]{\gdef\grothFirst{#1}\gdef\grothLast{#2}}
\newcommand{\foldertitle}[1]{\gdef\grothFoldertitle{#1}}
\gdef\grothFolder{?}\gdef\grothBatch{?}\gdef\grothFirst{?}\gdef\grothLast{?}\gdef\grothFoldertitle{}

% --- Le résumé --------------------------------------------------------------
%
% Il ouvre la lecture modernisée, sous le titre « Résumé », et il est composé
% en petit. Ce n'est pas un ornement : c'est le seul passage du document qui
% ne soit pas des mathématiques, et un lecteur doit pouvoir le reconnaître
% avant de l'avoir lu — puis le sauter s'il connaît déjà le sujet, ou s'y
% arrêter s'il ne le connaît pas. Le filet qui le suit marque la reprise du
% corps.

\newenvironment{resume}
  {\small\setlength{\parskip}{0.45em}}
  {\par\medskip\hrule\medskip}

% --- Mots-clés ---------------------------------------------------------------
%
% En anglais, en clôture du résumé de la lecture modernisée : le vocabulaire
% moderne sous lequel chercher ce que les feuillets construisent. Le manifeste
% du site (`npm run manifest`) les extrait pour étiqueter la cote entière —
% cette ligne est la source unique des tags, il n'y a pas de fichier à côté.
\newcommand{\keywords}[1]{\par\smallskip\noindent
  {\footnotesize\textsc{Keywords} --- \textit{#1}}\par}

% --- L'appareil critique ----------------------------------------------------

% Le numéro de feuillet, en marge. C'est l'ancre qui permet de régler un
% désaccord sur une formule en regardant *un* feuillet plutôt qu'une chemise
% de six cents. Le site s'en sert aussi pour tourner les pages du fac-similé
% quand on fait défiler la transcription.
\newcommand{\leaf}[1]{%
  \marginnote{\footnotesize\color{gray!70}\textsf{#1}}%
  \label{leaf:#1}\ignorespaces}

% Illisible : on ne devine pas. Un mot inventé qui se lit comme les autres est
% le pire résultat possible.
\newcommand{\ill}{\textcolor{red!60!black}{[\,\ldots\,]}}

% Lecture douteuse : le mot est donné, et signalé comme incertain.
\newcommand{\uncertain}[1]{\uline{#1}}

% Ajout de l'éditeur — une lettre manquante, un mot sous-entendu.
\newcommand{\add}[1]{\textcolor{blue!50!black}{[#1]}}

% Biffé par Grothendieck lui-même. Ce qu'il a rayé fait partie du document :
% une rature dit souvent où la pensée a changé de direction.
\newcommand{\struck}[1]{\sout{#1}}

% Note du transcripteur, sur l'état matériel du feuillet.
\newcommand{\note}[1]{\par\smallskip{\small\color{gray!60!black}\textit{[#1]}}\par\smallskip}

% Note marginale de Grothendieck, distincte du corps du texte.
\newcommand{\marginal}[1]{\par\smallskip\noindent\hspace{1em}%
  {\small\color{gray!60!black}#1}\par\smallskip}

% --- Titre ------------------------------------------------------------------

\AtBeginDocument{%
  \begin{center}
    {\large\bfseries Fonds Alexandre Grothendieck --- cote n\textsuperscript{o}~\grothFolder}\\[2pt]
    {\itshape\grothFoldertitle}\\[4pt]
    {\small feuillets \grothFirst--\grothLast\ (lot \grothBatch)}\\[2pt]
    {\footnotesize\color{gray!60!black}Transcription établie d'après le fac-similé numérisé
     par l'Université de Montpellier.\\
     Les lectures incertaines sont soulignées, les illisibles notées [\,\ldots\,],
     les ajouts entre crochets.}
  \end{center}
  \vspace{1em}}

\endinput


\folder{128}
\batch{2}
\pages{21}{31}
\dating{s.d. — le groupe « Descentes » (126 à 133) est daté [avant 1970]}
\watermark{Édition de démonstration}
\foldertitle{Descente non plate. Modules formels. Cf Notes Murre-Levelt : notes manuscrites (s.d.)}

\begin{document}

\section*{Descente le long des voisinages infinitésimaux (suite)}
\note{titre de l'éditeur, qui ne figure pas sur les feuillets ; le lemme 1
et le corollaire 1 auxquels renvoient les pages 21 et 22 précèdent ce
lot et n'y figurent pas}

\page{21}

\emph{Corollaire 2.} Sous les conditions du lemme~1/, soit
$n \geqslant 1$ un entier, $\mathbb{M}_n$ le faisceau « matrices carrées de
degré $n$ à coeff. dans \ill{} », $\mathrm{Gl}_n = \mathbb{M}_n^{*}$, alors
on a une bijection \uncertain{respectant} l'unité
\[
  \mathrm{Ker}\bigl(H^1(A'/A, \mathrm{Gl}_n) \to H^1(A'_0/A_0, \mathrm{Gl}_n)\bigr)
  \to
  \mathrm{Ker}\bigl(H^1(A'/A, \mathbb{M}_n) \to H^1(A'_0/A_0, \mathbb{M}_n)\bigr) .
\]
Par suite, si \struck{$H^1(A'/A, \mathrm{Gl}_n) = e$, alors}
$H^1(A'_0/A_0, \mathrm{Gl}_n) = \{e\}$,
$H^1(A'_0/A_0, \mathbb{M}_n) = \{e\}$, on a une bijection
\[
  H^1(A'/A, \mathrm{Gl}_n) \xrightarrow{\ \sim\ } H^1(A'/A, \mathbb{M}_n)
  = \mathbb{M}_n\bigl(H^1(A'/A)\bigr)
\]
donc pour que $H^1(A'/A, \mathrm{Gl}_n) = \{e\}$, il f.\ et s.\ que
$H^1(A'/A, G_a) = \{e\}$.
\note{sous « $H^1(A'_0/A_0, \mathbb{M}_n) = \{e\}$ » une accolade et une
double flèche verticale renvoient à « $H^1(A'_0/A_0) = 0$ » ; devant
$\mathbb{M}_n(H^1(A'/A))$ un symbole biffé, illisible}

\emph{Démonstration.} Les deux membres sont égaux respectivement à :
\[
  H^1\bigl(Z C^{\cdot}(A'/A, \mathbb{M}_n)\bigr) / \mathrm{Gl}(n, C_0)
\]
et à :
\[
  H^1\bigl(Z C^{\cdot}(A'/A, \mathbb{M}_n)\bigr) / \mathbb{M}_n(C_0)
\]
\note{le symbole lu $Z$ devant $C^{\cdot}$ est douteux dans les trois
occurrences. Dans la preuve du lemme 1 (pages 19--20, lot 1), le même complexe
s'écrit nettement $I\,C^{\cdot}(A'/A)$, l'idéal $I$ fois le complexe de
cochaînes : la lecture $I$ est donc probable ici aussi, à confronter aux pages}
et les opérations de $\mathrm{Gl}(n, C_0)$ sur
$H^1(Z C^{\cdot}(A'/A, \mathbb{M}_n))$ sont induites par celles de
$\mathbb{M}_n(C_0)$. De plus un élément de $\mathbb{M}_n(C_0)$ provenant de
$\mathbb{M}_n(B)$ opère trivialement. Donc il suffit \ill{} \ill{} \ill{}
\add{en} utilisant l'hypothèse \struck{\ill{} que}
\note{« utilisant l'hypothèse » en interligne ; la phrase s'arrête sur le
trait qui sépare la note de bas de page}

\marginal{/ On suppose que
$\mathbb{M}_n(C_0) = \mathbb{M}_n(C_0)^{*} + \mathrm{Im}\bigl(\mathbb{M}_n(B) \to \mathbb{M}_n(C_0)\bigr)$
[ce qui est vérifié sous les conditions signalées en N.B.\ au lemme~1]}
\note{note de bas de page de sa main, sous un trait, appelée par la barre
oblique placée après « lemme~1 » dans l'énoncé ; devant « Im » un mot
noirci}

\page{22}

\emph{Corollaire 3.} Soit $N$ un ens.\ d'indices, \struck{et} supposons
$E = A^{(N)}$, et supposons que l'on ait
\[
  \mathrm{End}_{C_0}(E_{C_0}) = \mathrm{Aut}_{C_0}(E_{C_0})
  + \mathrm{Im}\bigl(\mathrm{End}_{B_0}(E_{B_0}) \to \mathrm{End}_{C_0}(E_{C_0})\bigr) .
\]
Sous ces conditions, on a une bijection can., respectant les éléments
unités,
\[
  \mathrm{Ker}\bigl(H^1(A'/A, \underline{\mathrm{Aut}}(E)) \to H^1(A'_0/A_0, \underline{\mathrm{Aut}}(E))\bigr)
  \to
\]
\[
  \mathrm{Ker}\bigl(H^1(A'/A, \underline{\mathrm{End}}(E)) \to H^1(A'_0/A_0, \underline{\mathrm{End}}(E))\bigr)
\]
\[
  \simeq \Bigl[\bigl(\mathrm{Ker}(H^1(A'/A) \to H^1(A'_0/A_0))\bigr)^{(N)}\Bigr]^{N}
\]
\note{le $\simeq$ est écrit verticalement sous le second noyau ; les
exposants $(N)$ et $N$ sont repassés à l'encre foncée}

[Donc le noyau du premier nous est connu [i.e.\ les données de descente
sur $E_{A'}$ qui sont \uncertain{triviales} sur $E_{A'_0}$
\struck{\ill{}} \add{et} donnant un Module libre
\struck{$\simeq E_{C_0}$} sur $C_0$], ses éléments sont effectives et
donnent un Module libre $\simeq E_B$ sur $B$] \uncertain{lesquels} \ill{}
\ldots\ $\mathrm{Ker}(H^1(A'/A) \to H^1(A'_0/A_0)) = 0$]
\note{« telles que sur $E_{A'_0}$ » (lecture douteuse du premier mot) et
un mot au-dessus de « effectives », biffé, sont en interligne ; les
crochets ne se referment pas tous nettement}

En particulier, si $H^1(A'_0/A_0, \underline{\mathrm{Aut}}(E)) = \{e\}$ et
$H^1(A'_0/A_0) = 0$, on a une bijection
\[
  H^1(A'/A, \underline{\mathrm{Aut}}(E)) \simeq H^1(A'/A, \underline{\mathrm{End}}(E))
  \simeq \Bigl[H^1(A'/A)^{(N)}\Bigr]^{N} .
\]
\note{dans $H^1(A'/A, \underline{\mathrm{End}}(E))$ les indices $0$ de
$A'_0/A_0$ semblent effacés par des pâtés d'encre ; le second $\simeq$ est
vertical}

\emph{Démonstration.} Même principe que précédemment, en utilisant les
suites exactes de groupes simpliciaux

\begin{tikzcd}[column sep=small, row sep=small, nodes={font=\scriptsize}]
e \arrow[r] & C^{\cdot}(A'/A, \underline{\mathrm{Hom}}(E, IE)) \arrow[r] \arrow[d] & C^{\cdot}(A'/A, \underline{\mathrm{Aut}}(E)) \arrow[r] \arrow[d] & C^{\cdot}(A'_0/A_0, \underline{\mathrm{Aut}}(E)) \arrow[r] & e \\
0 \arrow[r] & C^{\cdot}(A'/A, \underline{\mathrm{Hom}}(E, IE)) \arrow[r] & C^{\cdot}(A'/A, \underline{\mathrm{End}}(E)) \arrow[r] & C^{\cdot}(A'_0/A_0, \underline{\mathrm{End}}(E)) \arrow[r] & 0
\end{tikzcd}

\note{la flèche verticale de droite n'est pas dessinée sur la page ; les
deux autres le sont}

\page{23}

\emph{Lemme 2.} Soit $B$ un anneau ½ local, $C$ une algèbre finie sur $B$,
telle que $C$ soit ½ local, et les idéaux maximaux \struck{\ill{}}
$\mathfrak{n}$ de $C$ sont induits par les idéaux maximaux \struck{\ill{}}
$\mathfrak{m}$ de $B$. Supposons les conditions suivantes satisfaites :
\begin{itemize}
  \item[(i)] Les corps $k(\mathfrak{m})$ sont infinis.
  \item[(ii)] Deux idéaux maximaux $\mathfrak{n}, \mathfrak{n}'$ distincts de
  $C$ induisent des idéaux maximaux $\mathfrak{m}, \mathfrak{m}'$ distincts de
  $B$, et pour tout $\mathfrak{n}$, l'extension résiduelle
  $k(\mathfrak{n})/k(\mathfrak{m})$ est radicielle (p.ex.\ $C$ fini et
  radiciel sur $B$ ½ local).
\end{itemize}
Alors pour tout $n \geqslant 1$, et tout élément $u \in \mathbb{M}_n(C)$,
$\exists\, v \in \mathbb{M}_n(B)$ tel que $u + \varphi(v) \in \mathbb{M}_n(C)^{*}$.
\note{« $\mathfrak{n}$ » et « $\mathfrak{m}$ » sont écrits en interligne
au-dessus de deux mots biffés}

\emph{Dém.} Soient $\mathfrak{m}_i$ \uncertain{(resp.\ $\mathfrak{n}_i$)}
les idéaux maximaux de $B$ \ill{} \ill{} \ill{} de $C$. \ill{} \ill{}
\ill{} $\to$ fini ; \ill{} \ill{} donc par th.\ \ill{} \ill{} \ill{}
d'éléments $v_i \in \mathbb{M}_n(k(\mathfrak{m}_i))$, il existe un
$v \in \mathbb{M}_n(B)$ qui les donne. D'ailleurs pour $u' \in
\mathbb{M}_n(C)$, le fait que ce soit \uncertain{inversible} ne dépend que
de l'inversibilité de ses images dans $\mathbb{M}_n(C/\mathrm{rad}\,C)$ !
Cela nous ramène au cas où $C$ est un \uncertain{produit} de corps, et $B$
aussi. \uncertain{On} \ill{} \ill{} au cas $B$ un corps, $B = k$.

\page{24}

\ill{} --- $v \in \mathrm{End}_C(E_C)$

Il existe $u \in \mathrm{End}_B(E)$ tel que $v + \struck{\ill{}}\, u_C$
soit \emph{inversible}.

\emph{Dém.} Si $E \simeq B^{(N)}$, \ill{} \ill{}
$\mathrm{End}_B(E) \simeq E^{N}$, donc la \uncertain{remarque} \ill{}
\uncertain{l'élément} $u \in \underline{\mathrm{End}}(E)(B)$ \ill{}
\ill{}
$\underline{\mathrm{End}}(E)(k(\mathfrak{m}_i))$ \ill{} \ill{}
[\,\ill{} \ill{} \ill{} : \ill{} \ill{} \ill{}
$\alpha_i \in E(k(\mathfrak{m}_i))$, et \ill{} \ill{} \ill{}, \ill{} :
$E \simeq B^{(I)}$, \ill{} \ill{} \ill{} \ill{} \ill{} dans
$k(\mathfrak{m}_i)$ \ill{} $\Rightarrow$\,]. \ill{} \ill{}, \ill{} ?
\ill{} l'hypothèse \ill{} \ill{} \ill{}$(C)$ \ill{}, \ill{}
\ill{} \ill{} \ill{} \ill{} \ill{}$(C)$. Car \ill{} \ill{} \ill{}
\ill{} \ill{} \ill{} $C$ \ill{} $B$, \ill{} \ill{} \ill{} \ill{}
\struck{\ill{}}, \ill{} : l'hypothèse \ill{}, \ill{} \ill{} \ill{}
\ill{} \ill{} $u = \mathrm{id}_E - v$.
\struck{\ill{} \ill{} \ill{} $B \rightrightarrows C$, \ill{} $B$ \ill{}
\ill{} ; \ill{} (i), \ill{} \ill{} $u = \lambda\, \mathrm{id}_E$ \ill{}}
\note{page très rapide ; trois lignes à la fin de la démonstration sont
biffées d'un trait ; au-dessus de l'une d'elles un $C$ en interligne}

N.B. Il faudrait \uncertain{dégager} la possibilité d'une extension
résiduelle transcendante $k(t)$, qui \uncertain{fera} \ill{} \ill{}
\ill{} en (i) \ldots

\page{25}

\emph{Théorème.} $A$ anneau \struck{artinien} ½ local noethérien,
\struck{\ill{}} $\mathfrak{m}$ son radical, $A'$ une algèbre finie sur $A$.
Pour tout entier $n$, on pose $A_n = A/\mathfrak{m}^{n+1}$,
$A'_n = A'/\mathfrak{m}^{n+1}A' = A' \otimes_A A_n$.
\note{« ½ local noethérien » est écrit au-dessus d'« artinien », biffé}

[a) Soit $m$ un entier $\geqslant 1$. Les conditions suivantes sont
équivalentes
\begin{itemize}
  \item[(i)] $\forall n$, on a $H^1(A'_n/A_n) = 0$
  \item[(ii$_1$)] $\forall n$, on a $H^1(A'_n/A_n, \mathrm{Gl}_1) = 0$
  \item[(ii$_m$)] $\forall n$ on a $H^1(A'_n/A_n, \mathrm{Gl}(m)) = \{e\}$
  \item[(iii)] $\forall n$, désignant par $B^{(n)}$ l'anneau
  $H^0(A'_n/A_n)$, le morphisme $\operatorname{Spec} A'_n \to
  \operatorname{Spec} B^{(n)}$ est un morphisme de descente effective pour
  la catégorie fibrée des Modules \uncertain{qu.\ cohér.}\,/\,plats.
\end{itemize}
\note{un crochet vertical embrasse les deux lignes « a) Soit $m$\ldots » et
« $A'\otimes_A A_n$ » ; « qu.\ cohér. » et « plats » sont écrits l'un
au-dessus de l'autre, séparés par un trait, comme une fraction}

b) Ces conditions impliquent
\begin{itemize}
  \item[(i)] $H^1(A'/A) = 0$
  \item[(ii)] $H^1(A'/A, \mathrm{Gl}_m) = 0$, plus généralement pour tout
  entier $m$ $H^1(A'/A, \mathrm{Gl}(m)) = \{e\}$
  \item[(iii)] Désignant par $B$ l'anneau $H^0(A'/A)$, le morphisme
  $\operatorname{Spec}(A') \to \operatorname{Spec}(B)$ est un morphisme de
  descente effective pour la catégorie des modules plats \emph{de type
  fini}.
\end{itemize}
\note{« de type fini » est souligné et entouré au crayon}

c) Supposons $A$ artinien, et les conditions de a) satisfaites. Alors pour
tout $B$-\uncertain{préschéma} $X$ plat sur $B$, posant
$X' = X \otimes_B A'$, le morphisme $X' \to X$ est un morphisme de descente
effective pour la catégorie fibrée des Modules plats quasi-cohérents.
\note{« artinien » est entouré au crayon}

\marginal{N.B. Dans a) les \uncertain{hypothèses} \ill{} \ill{} \ill{}
\ill{} \ill{} \ill{} \ill{} \ill{} \ill{} [$A$ anneau \struck{\ill{}}
\ill{} \uncertain{local noethérien} \uncertain{quelconque}, $A'$ \ill{}
$A$ \ldots]}
\note{écrit en diagonale dans la marge gauche, de bas en haut, face à
(ii)--(iii) de a) ; un mot souligné d'un double trait. Plus bas, dans la
même marge et en diagonale, un mot \ill{} qui semble relié à la ligne
« morphisme de descente » de (iii)}

\page{26}

(a)\ \emph{Démonstration.} Par les lemmes 1 et 2, l'équivalence
des conditions (i) \struck{(ii)} (ii$_m$) \struck{(iii)} est immédiate par
récurrence sur $n$, compte tenu que (i) elle est vraie pour $n = 0$
(\emph{descente fid.\ plate}) (ii) les conditions
\[
  \mathbb{M}_m(B^{(n)}) = \mathbb{M}_m(B^{(n)})^{*} + \mathrm{Im}\,\mathbb{M}_m(B^{(n)})
\]
\struck{relatif à} sont satisfaites pour tout $n$, en vertu du lemme~2 (ii)
\ill{} \ill{} et de la théorie des anneaux artiniens \ldots
\note{au-dessus de « et 2 », « cor.~2 » en interligne, relié par une
accolade ; la seconde égalité est écrite de façon très serrée et son
dernier terme est douteux}

La condition (ii$_m$) signifie que $\forall n$,
$\operatorname{Spec} A'_n \to \operatorname{Spec} B^{(n)}$ est un
morphisme de descente effective \struck{\ill{}} pour la catégorie des
Modules \uncertain{qu.\ coh.}\ \emph{libres} de \struck{type fini}
\uncertain{rang fini}. Cela signifie aussi que tout $A'_n$-module
\uncertain{$A'^{\,m}_n$} muni d'une donnée de descente relative à
$A'_n/B^{(n)}$, est effectif et donne un module \ill{} $E_{B^{(n)}}$ sur
$B^{(n)}$ (\uncertain{qui} \ill{} \ill{} après changement de base
$B^{(n)} \to A'_n$). \ill{} \ill{} \ill{} \ill{} \ill{} \ill{} \ill{}
\note{« rang fini » en interligne au-dessus de « type fini », biffé ; la
parenthèse est en interligne, reliée par un trait}

Je dis que les modules descendus \ill{} \ill{} \ill{} \ill{} \ill{},
\ill{} \ill{} $\mathrm{Ker}(B^{(n)})$ l'est, \ill{} \ill{} \ill{}
\ill{} \ill{} \ill{} \ill{} \ill{}, \ill{} corps, \ill{} \ill{}
\ill{} \ill{} \ill{} \ldots
\marginal{\ill{} : \ill{}}
\note{en marge gauche, séparés du texte par un trait vertical, deux mots
de sa main ; « de corps, puis » en interligne au-dessus d'un mot}

Donc il est clair que (iii) $\Longrightarrow$ (ii$_m$). Pour voir que

\page{27}

\struck{(iii)} (i) $\Longrightarrow$ (iii) pour \ill{} \ill{} donné ;
[\,\ill{} \ill{} que (i) est stable par changement de base, et (iii) est
\emph{invariant} par changement \struck{plat} \emph{fid.}\ plat de la base.
Cela nous ramène au cas où les extensions résiduelles de $A'/A$ sont
triviales. \struck{Mais \ill{}}
$A = A_n = \underline{B}^{(n)}$, et
\uncertain{on} peut supposer [$A$ local [en le remplaçant par ses
localisés]. Alors on sait que $E_{A'_n}$ \struck{plat sur} $A'_n$ est muni
d'une donnée de descente, \struck{\ill{}} relative à $A'_n/A_n$ [i.e.\ :
$A'_n/B^{(n)}$]
\note{un crochet ouvert devant « \ill{} \ill{} que (i) » ; « fid. » en
interligne au-dessus de « plat », biffé ; « $A = A_n = B^{(n)}$, et » en
interligne, relié par un trait}

\ill{} $A'_n$-module \emph{libre} [car il est \ill{} \ill{} \ill{}
\ill{} sur un anneau local \ill{} \ill{} \ill{} défini par \ill{} \ill{}
$\ill{}$, on peut \ill{} ramener au cas \ill{} la \uncertain{spectre}
d'un corps, \ill{} \ill{} \ill{} \ill{} \ill{} \ill{} \ill{} de
descente.]

Je dis que \uncertain{ces} effectivités par $A'_n/B^{(n)}$ ; donc
\uncertain{celles} \ill{} \ill{} \ill{} \ill{} \ill{} sur $B^{(n)}$ :
cela résulte des lemme~1, cor.~3, compte tenu du lemme~2.

\uncertain{On} a ainsi prouvé a).

(b)\ Pour prouver (i), \uncertain{prenons} la formation de
$H^1(A'/A)$ commutant : l'extension plate de la base, on peut supposer $A$
\emph{complet}. Alors \struck{\ill{}}

\page{28}

\[
  C^{\cdot}(A'/A) = \varprojlim_n C^{\cdot}(A'_n/A_n) ,
\]
donc en vertu de M.L.\ on a
\[
  H^1(A'/A) = H^1\bigl(C^{\cdot}(A'/A)\bigr)
  = \varprojlim H^1\bigl(C^{\cdot}(A'_n/A_n)\bigr)
  = \varprojlim_n H^1(A'_n/A_n) = 0 .
\]

On peut aussi \uncertain{déduire} formellement (i) de (iii) sans difficulté
(en faisant $m = 1$, et $m = 2$).

(ii) \uncertain{Vm} résulte de (iii), par le même argument que dans a).

(iii) \struck{\ill{}} On voit comme dans (i) que l'on peut supposer $A$
\struck{\emph{local}} \emph{complet}, \struck{et évidemment}
\struck{Enfin, on} On peut aussi supposer $A = B$, et $A$ \emph{local}. Soit
$E' = E_{A'}$ plat de type fini sur $A'$, muni d'une donnée de descente, et
soit
\[
  E = H^0_D(A'/A, E') .
\]
Il faut prouver 1°) $E$ plat sur $A$ 2°) $E \otimes_A A' \to E'$ est un
isomorphisme. Or on a
\[
  C^{\cdot}_D(A'/A, E') \simeq \varprojlim_n C^{\cdot}_D(A'_n/A_n, E'_n)
\]
d'où par M.L.
\[
  E = H^0_D(A'/A, E') = \varprojlim_n H^0(A'_n/A_n, E'_n) .
\]
\note{dans « $E' = E_{A'}$ », l'indice de $E'$ est noirci}

Posons
\[
  E^{(n)} = H^0(A'_n/A_n, E'_n)
\]
alors grâce à l'hypothèse a), sous la forme (iii), on sait que $E^{(n)}$
est un $B^{(n)}$-module plat de t.f., et
$E^{(n)} \otimes_{B^{(n)}} A'_n \to E'_n$ est un isom.

\page{29}

Moyennant un lemme facile de passage à la limite, on en déduit que
$E = \varprojlim E^{(n)}$ est plat sur $A = \varprojlim_n B^{(n)}$.
\uncertain{Pour} voir enfin que $E \otimes_A A' \to E'$ est un
isomorphisme, on note que les deux membres sont \uncertain{respectivement}
les $\varprojlim$ des $E^{(n)} \otimes_{B^{(n)}} A'_n$ et des
$E'_n = E' \otimes_A A'_n$, et l'homomorphisme est la $\varprojlim$ des
\uncertain{isomorphismes} $E^{(n)} \otimes_{B^{(n)}} A'_n \to E'_n$, donc
est un \emph{isomorphisme}.

(c)\ \uncertain{Montrons} que $X' \to X$ est un morphisme de
descente pour la catégorie en question, \uncertain{puis} que c'est un
morphisme de descente effective, \ill{} \ill{} [De plus, on peut supposer
évidemment $A = B$ \ill{}, \ill{} \ill{} \ill{} \ill{} \ill{} pour les
$A_n$ \ldots, \ill{} l'hypothèse sous la forme a) (iii) \ldots\,]

\[
  A \longrightarrow A' \rightrightarrows A''
\]
\note{ligne isolée au bas de la page, sans texte}

\page{31}

\note{feuillet d'une autre main, écrit proprement en anglais, sans
intervention de Grothendieck repérée ; vu le titre du dossier, il s'agit
vraisemblablement d'une page des notes Murre--Levelt, mais rien sur la
feuille n'en nomme l'auteur. Pour $A' = k + V = k[X_1, \ldots, X_n]/(X_i)^2$
et $I$ un « equivalence ideal » de $A' \otimes_k A'$, engendré par des
$F_\alpha(X, Y) = L'_\alpha(X) + L''_\alpha(Y) + Q_\alpha(X, Y)$ (parties
linéaires et bilinéaire), la condition $p_{13}(I) \subset p_{12}(I) +
p_{23}(I)$ fournit des $G_{\alpha\beta}$, $H_{\alpha\gamma}$ tels que
$F_\alpha(X, Z) \equiv \sum_\beta F_\beta(X, Y) G_{\alpha\beta}(X, Y, Z) +
\sum_\gamma F_\gamma(Y, Z) H_{\alpha\gamma}(X, Y, Z)$ modulo les carrés ;
on en tire une expression de $Q_\alpha(X, Z)$ par des formes linéaires
$g_{\alpha\beta}$, $h_{\alpha\gamma}$, on remplace les générateurs par
$F'_\alpha = F_\alpha - \sum_\beta F_\beta\, g_{\alpha\beta} - \ldots$,
qui engendrent encore $I$, et, si $I \subset \Delta$, on conclut que
$F'_\alpha(X, Y) = L'_\alpha(X) - L'_\alpha(Y)$}

\end{document}
